% Lección: límite de una función en un punto (castellano).
%
% Sin preámbulo, sin \documentclass, sin \begin{document}: solo el contenido.
% La lección no sabe a qué salida va. Lo que no va envuelto sale en todas;
% \onlynotes{...} solo en los apuntes; `teaching` solo en las copias del
% profesor. \begin{frame} es una diapositiva en las diapositivas y un bloque
% de texto seguido en los apuntes.

\begin{frame}
\didactatitle{Límite de una función en un punto}

\bymedium{%
  $f(x)$ se acerca a $L$ cuando $x$ se acerca a $a$. La definición dice qué
  quiere decir «se acerca».
}{%
  La idea intuitiva es que $f(x)$ se acerca tanto como se quiera a $L$ cuando
  $x$ se acerca lo suficiente a $a$. La definición convierte esa frase en una
  desigualdad que se puede comprobar.
}

\begin{definition}[Límite]
Sea $f$ una función definida en un intervalo abierto que contiene a $a$,
salvo quizá en el propio $a$. Decimos que $f$ tiene \keyterm{límite} $L$ en
$a$, y escribimos $\lim_{x \to a} f(x) = L$, si para cada $\varepsilon > 0$
existe un $\delta > 0$ tal que
\begin{keyformula}
  0 < |x - a| < \delta \implies |f(x) - L| < \varepsilon .
\end{keyformula}
\end{definition}

% Dos párrafos: para eso está la forma de entorno. \onlynotes{...} es para
% una frase o dos sin líneas en blanco dentro.
\begin{notesonly}
La condición $0 < |x - a|$ deja fuera el propio punto $a$: el límite no
depende de lo que valga $f$ en $a$, ni siquiera de que esté definida allí.
Solo mira lo que pasa \emph{cerca} de $a$. Por eso tiene sentido preguntarse
por el límite de $(x^2 - 1)/(x - 1)$ en $x = 1$, donde el cociente no está
definido.

Además, el límite, si existe, es único: dos valores distintos $L \neq L'$ no
pueden estar a la vez a menos de $\varepsilon = |L - L'|/2$ de $f(x)$.
\end{notesonly}

\begin{teaching}[Ritmo]
Veinte minutos. Dibuja la banda $(L - \varepsilon, L + \varepsilon)$ antes de
escribir la fórmula: el orden de los cuantificadores (primero te dan
$\varepsilon$, luego eliges $\delta$) se entiende mejor en la pizarra.
\end{teaching}
\end{frame}
